\chapter{Etat de l'art : projets similaires}

Il convient de rappeler qu'un projet de base de données au sein de la BnF s'inscrit dans la continuité de travaux similaires réalisés auparavant. De même, l'utilisation d'\emph{Heurist} n'est pas inédite dans le domaine des sciences humaines et sociales, notamment à la BnF. 


\section{Les bases de données à la BnF}

Une base de données est un ensemble structuré de fichiers regroupant des informations ayant certains caractères en commun. Cet ensemble représente un système d'informations sélectionnées de façon à pouvoir être consultées par des utilisateurs ou des programmes\footcite{BDD_def}. Du fait des activités et de l'envergure de la BnF, la notion de base de données est très large dans le contexte de cette institution, et englobe toutes les données structurées, y compris les ressources les plus communément utilisées, telles que le catalogue général\index{Catalogue général de la BnF} ou Gallica\index{Gallica}.  
Ainsi, la base de données dont mon stage fait l'objet s'inscrit plutôt dans le contexte de la structuration de données en lien avec un projet de recherche. En effet, plusieurs programmes de recherche de la BnF ont abouti à la création de bases de données, qui donnent accès à des ressources spécialisées. C'est le cas de Mandragore, base de données scientifique dédiée à l’identification et à l'indexation de décors de manuscrits du département des Manuscrits et de la Bibliothèque de l’Arsenal. La base est mise en place en 1989 par le centre de recherche sur le manuscrit enluminé (CRME), qui est le fruit d'une collaboration entre le département des Manuscrits de la BnF et le CNRS. 
Parmi les bases de données venant de programmes de recherche réalisés à la BnF, certaines couvrent des thématiques assez proches de celles du projet Heurist sur les factures Rondel. Elles ont pu servir d'inspiration pour ce projet, qui s'inscrit dans leur continuité. La base \enquote{reliures}, par exemple, recense les reliures françaises à décor du XVI\textsuperscript{e} siècle au début du XIX\textsuperscript{e} siècle conservées à la réserve des livres rares de la BnF. Cette base a été utile à ma réflexion, car elle traite également du livre ancien dans sa matérialité et dans sa circulation, et s'intéresse aussi aux personnes exerçant des métiers liés au livre ancien. Dans la même veine, la base BP16 (\enquote{Bibliographie des éditions parisiennes du 16ème siècle}) recense, à partir des travaux de Philippe Renouard, toute la production imprimée à Paris au XVI\textsuperscript{e} siècle. Cette base a aussi été une source d'inspiration : je me suis inspiré de son exemple pour accorder une place centrale, dans ma base, aux exemplaires plus qu'aux oeuvres elles-mêmes. 
De plus, la base BP16 accorde une certaine attention aux libraires-éditeurs parisiens du XVI\textsuperscript{e} siècle, et l'un des objectifs de l'exploitation des factures Rondel est de faire une cartographie des libraires, notamment parisiens, fournisseurs de Rondel, ce qui place les deux projets dans la même continuité historique, et dans la continuité des pratiques professionnelles de l'institution. 

\section{Présentation et fonctionnement d'\emph{Heurist}}

\subsection{Histoire et présentation}


\emph{Heurist}\index{Heurist}, logiciel utilisé pour réaliser ce projet, est un système de gestion de base de données (SGBD) créé en 2005 par Ian Johnson, professeur à l'université de Sydney. \emph{Heurist} permet l'élaboration de bases de données relationnelles MySQL richement structurées\footcite{huma-num_heurist_2026}. Une base de données relationnelles (BDDR) est un type de base de données, qui se distingue par sa capacité à mettre en relation différentes données au sein de cette dernière.\footcite{marina_hervieu_2025}. Théorisé par un article intitulé \enquote{A Relational Model of Data for Large Shared Data Banks} et écrit par Edgar Codd en 1970, le modèle relationnel entend structurer les données sous forme de domaines (ensemble de valeurs) et de relations (tables). Ce modèle ne s'est pas imposé de manière universelle, mais a fortement influé sur les outils et pratiques de la base de données, du fait de ses nombreux atouts. En effet, la BDDR offre une représentation simple des données, n'étant constituée que de domaines et de relations, qui reste rigoureuse mathématiquement, puisqu'elle s'appuie sur la théorie des ensembles et des relations. Ce modèle est aussi révolutionnaire par sa capacité à décrire et à manipuler facilement les données, grâce à un langage unique qui s'est progressivement imposé\footcite[p. 45-47]{Hainaut_2015}. Il s'agit du SQL (\emph{Structured Query Language})\footcite[p. 191]{Hainaut_2015}, utilisé dès la fin des années 1970 et permettant de déclarer, de manipuler et d'interroger les données, mais aussi de contrôler leur accès\footcite[p. 166]{Soutou_2022}. 
Ces bases de données s'organisent sous la forme d'un ou plusieurs tableaux, appelés \enquote{tables}, contenant un ensemble de lignes. Une ligne est une suite de valeurs, et regroupe des informations concernant un objet. L'ensemble des valeurs de même type constituent une colonne. Les tables peuvent être liées entre elles par des colonnes particulières, appelées \enquote{clés étrangères}, dont le rôle est de référencer une table dans une autre\footcite[p. 44-54]{Hainaut_2015}. 


\emph{Heurist}\index{Heurist} est utilisé dans notre projet car il s'agit d'un logiciel particulièrement adapté aux projets de recherche en sciences humaines et sociales, et cela à plusieurs titres. En effet, ce logiciel a pour but de permettre aux chercheuses et aux chercheurs en sciences sociales de gérer leurs données acquises sur le terrain. L'enjeu est de leur redonner le contrôle sur leurs données, au lieu de le déléguer aux personnes en charge du développement informatique. Cela est permis par sa relative facilité d'utilisation : \emph{Heurist} est accessible par simple navigateur web et ne requiert aucune installation. De plus, il ne demande pas de compétences en programmation, c'est pourquoi il est accessible à un public relativement  large\footcite{huma-num_heurist_2026}. Ce logiciel est également apprécié du fait de son caractère \emph{Open source}, et de sa compatibilité avec les objectifs de la science ouverte. En effet, ses données peuvent être exportées dans des formats ouverts, comme le CSV, et sont accessibles sur le web et interopérables, notamment par la possibilité de publier un site web à partir d'une base de données\footcite{paillusson_introduction_2022}. 
\emph{Heurist} a été pensé pour résoudre des problématiques de recherche en SHS, c'est pourquoi il intègre des fonctionnalités en phase avec ces questions, telles que des visualisations temporelles ou spatiales\footcite{paillusson_introduction_2022}. 
Enfin, la considération suivante peut paraître secondaire, mais les logiciels spécifiques liés aux bases de données sont souvent appréciés par les chercheurs, car leurs formulaires de saisie sont plus \enquote{avenants}\footcite[Voir \enquote{Tableur ou gestionnaire de bases de données ?}]{Lemercier_Zalc_2008}. Ce détail a son importance, et fait partie des raisons pour lesquelles Heurist semble plus approprié qu'une base de données brute pour la recherche en SHS. 

\subsection{Les \emph{Records} et \emph{Record types}}

Les éléments au centre d'une base \emph{Heurist}\index{Heurist} sont les \emph{Records} et les \emph{Record types}. Le terme de \emph{Record type}, littéralement \enquote{type d'enregistrement}, est l'équivalent dans Heurist des tables d'une base de données, à savoir un ensemble d'objets de même nature. Ces Record types contiennent des \enquote{champs} (fields), c'est-à-dire des composantes qui fournissent leur structure, et correspondent aux colonnes d'une base de données. Les Records types et leurs champs sont d'ailleurs créés par la fonction \enquote{Design} du logiciel. Les champs sont aussi appelés \enquote{attributs} dans le contexte d'une base de données. Ces champs sont complétés par des \emph{Records}, qui sont des instances de ces tables, constitués de données qui remplissent les différents champs. Lorsqu'un champ est obligatoire (\enquote{mandatory}), l'enregistrement ne peut pas exister si ce champ là n'est pas rempli par des données. Le champ peut aussi être recommandé (\enquote{recommended}), facultatif (\enquote{optional}) ou caché (\enquote{hidden}). Dans le dernier cas, le champ n'apparaît pas lorsque l'on affiche les enregistrements de cette table. 

\subsection{Les types de données}
\label{types donnees}

Comme une base de données classique, \emph{Heurist}\index{Heurist} offre la possibilité de choisir le type de données que l'on veut utiliser pour chaque champ. Cependant, les types de données proposés par Heurist sont très spécifiques à ce logiciel, c'est pourquoi il convient de les détailler ici. Tout d'abord, \emph{Heurist}\index{Heurist} propose deux types de données textuelles : text (single-line), qui n'autorise qu'une seule ligne de texte, et sert à renseigner des noms ou des descriptions courtes, et memo text (multi-line), utile pour les paragraphes. Pour les nombres, le type de données \emph{Numeric} permet de renseigner des nombres positifs ou négatifs, entiers ou décimaux. En base de données relationnelles, les données numériques peuvent aussi être générées automatiquement dans une colonne dédiée, afin de donner à chaque enregistrement d'une table un identifiant unique. Il s'agit de générateurs de séquences de valeurs, appelés différemment selon le SGBD, mais souvent par le terme \emph{autoincrement}\footcite[p. 194]{Hainaut_2015}. \emph{Heurist}\index{Heurist} permet aussi de générer des identifiants numériques pour chaque enregistrement d'une table, en rendant un champ numérique obligatoire et en cochant la case \emph{Increment value by 1}. Il faut noter qu'un autre identifiant est généré automatiquement à chaque enregistrement dans \emph{Heurist}\index{Heurist}, peu importe son \emph{Record type} : l'\emph{Heurist identifier}, aussi appelé \emph{Heurist-ID}. Les \emph{Heurist-ID} servent surtout à être lus par le logiciel et à être reliés à des enregistrements déjà existants en cas d'imports de données en masse. La gestion des dates est très pratique dans \emph{Heurist}\index{Heurist}, grâce au type de données \enquote{Date/temporal}, qui permet d'avoir différentes échelles de précision sur les dates que l'on renseigne. En effet, on peut choisir si on entre une date exacte dans la base, au format YYYY-MM-DD, ou une date approximative, ce qui affichera dans la base le terme \enquote{circa} suivi de la date. Il est aussi possible d'indiquer une date avant ou après laquelle a eu lieu l'évènement qu'on souhaite dater, ce qui indiquera \enquote{before} ou \enquote{after} suivi de cette date. On peut aussi renseigner un intervalle de dates, et des degrés de certitude pour chaque date renseignée, à savoir \emph{Unknown} (inconnue), \emph{Attested} (attestée), \emph{Conjecture} (supposition) ou \emph{Measurement} (estimation). 
Les données de type \emph{Dropdown} sont l'une des spécificités d' \emph{Heurist}\index{Heurist}. Il s'agit d'un choix déroulant entre différents mots de vocabulaire, issus de listes que l'on définit dans une fonctionnalité spéciale du logiciel, appelée \emph{Vocabularies} (vocabulaires contrôlés). Les vocabulaires contrôlés permettent donc d'avoir des listes de mots, pré-conçues pour remplir certains champs, ce qui évite de tout devoir saisir à la main. Il existe également un type de données nommé \emph{File or remote URL} et permettant d'importer des fichiers ou une URL. 


\subsection{Les relations}

Certains types de données dans \emph{Heurist}\index{Heurist} servent à lier différents Record types entre eux, ce qui est l'intérêt d'une base de données relationnelle. Ces types de données sont le \emph{Record pointer} et le \emph{Relationship marker}. Les Record pointers sont l'équivalent dans \emph{Heurist}\index{Heurist} des clés étrangères dans une base de données classique : ils servent à lier une table à une autre. Lorsque l'on crée un champ de ce type, il faut d'abord choisir à quel \emph{Record type} on veut faire la liaison, et pour remplir ce champ, il faut sélectionner l'enregistrement auquel on veut lier celui qu'on est en train de créer. Les \emph{Relationship markers}, deuxième type de relation possible, consistent encore une fois à lier deux enregistrements entre eux, mais cette fois avec la possibilité de préciser quel type de relation existe entre ces enregistrements. Ce type de données permet également de décrire et de mettre des bornes chronologiques à ces relations. Tout comme les \emph{Dropdown}, les \emph{Relationship markers} font aussi appel à des vocabulaires contrôlés que l'utilisateur peut concevoir lui-même au préalable. 

En somme, \emph{Heurist} est un logiciel permettant de créer et de gérer des bases de données relationnelles, selon des types de données qui lui sont propres. 

\section{Projets \emph{Heurist} similaires à celui des factures Rondel}

Cette utilisation fréquente d'\emph{Heurist} dans des projets de recherche en SHS fait que le nôtre s'inscrit dans la continuité de projets précédents, traitant les mêmes thématiques. 

Le projet DEF-19 (Dictionnaire des éditeurs français du XIX\textsuperscript{e} siècle) a été financé par l'ANR entre septembre 2014 et septembre 2019, et a notamment consisté à la mise en place d'une base de données \emph{Heurist}\index{Heurist}. Celle-ci est composée de deux tables : une consacrée aux personnes (éditeurs) et une aux établissements (maisons d'édition). Ainsi, chaque établissement est lié aux personnes qui le composent. Cette base présente également les données sous forme cartographique, où l'utilisateur peut affiner sa recherche pour sélectionner ce qu'il souhaite voir sur la carte\footcite{noauthor_dictionnaire_nodate}. Cette base nous a également inspiré, car elle traite, une fois de plus, les mêmes thèmes que la nôtre. Nous nous sommes donc inspirés de sa structure, qui était assez claire, pour modéliser les liens entre les personnes et les établissements. 

\section{Projets \emph{Heurist} déjà réalisés à la BnF}


La conception d'une base \emph{Heurist} au sein du département des Arts du spectacle participe du déploiement du logiciel à la BnF. Au-delà des autres raisons évoquées précédemment, qui font d'\emph{Heurist} un logiciel adéquat pour les projets de recherche en sciences humaines, le logiciel convient parfaitement au contexte humain et professionnel d'une institution patrimoniale aussi vaste que la BnF. En effet, \emph{Heurist} peut rapidement être pris en main par différentes personnes, sans qu'elles aient nécessairement de connaissances techniques sur les bases de données. On peut donc aisément imaginer des membres d'un ou plusieurs services (ou département) utiliser la même base \emph{Heurist}, que ce soit pour l'implémenter ou pour chercher des informations dedans. \emph{Heurist} peut également être apprécié à la BnF pour sa capacité à s'imposer dans le temps. Son utilisation est très guidée, et les bases peuvent contenir une grande quantité de documentation. En effet, le logiciel prévoit, au sein de chaque élément de la base, (\emph{Record type}, \emph{fields}), des espaces de textes pouvant accueillir des indications sur l'utilisation qu'il faut en avoir. 
C'est pourquoi une base \emph{Heurist} pourrait perdurer et survivre par exemple à un renouvellement d'équipe. Le logiciel se prête donc très bien à l'exploitation d'un corpus aussi large que les factures Rondel, qui devra être réalisée sur une longue période et par plusieurs personnes. 
Ainsi, plusieurs bases \emph{Heurist}\index{Heurist} ont déjà été créées au sein de la BnF, avec différents niveaux d'aboutissement. 

\subsection{Collections de collectionneurs}

La base \emph{Heurist} \enquote{Collections de collectionneurs} est une composante du projet de recherche éponyme, mené par la bibliothèque de l'Arsenal et inscrit au plan quadriennal de la recherche pour la période 2020-2023\footcite{noauthor_plan_nodate}. 
Ce projet est assez proche, dans sa thématique, de celui entrepris pendant mon stage, et part du constat que les catalogues actuels de la BnF ne sont pas adaptés à l'analyse des anciennes collections bibliophiliques qu'elle conserve, alors qu'elles font l'objet de travaux par des chercheurs et des conservateurs. L'objectif est donc de créer un modèle de données spécifique pour permettre une meilleure description, visualisation et gestion de ces collections de bibliophiles. Le but est également de s'intégrer aux systèmes d'information de la BnF et d'assurer la pérennisation des données.

En effet, cette base \emph{Heurist}\index{Heurist} décrit avec précision les ouvrages issus de ces collections. Les informations sont requêtables grâce aux recherches filtrées permises par le logiciel, portant une attention particulière à la matérialité des livres (reliure, tranche, décors)\footcite{noauthor_collection_nodate}. 


\subsection{Base des trésors monétaires}

La base de données \enquote{Trouvailles monétaires} recense l'ensemble des dépôts monétaires découverts sur le territoire français et les données venant des archives du département des Monnaies, médailles et antiques de la BnF. Le département porte le programme de recherche éponyme depuis 1978, avec pour ambition d'étudier et de valoriser les découvertes numismatiques. La base s'est développée grâce au soutien de deux plans quadriennaux de la recherche entre 2016 et 2023, et se poursuit grâce à l'ajout de nouvelles données, notamment via le projet Archinumis (2024-2027), qui enrichit la base par l'étude d'archives numismatiques\footcite{noauthor_base_nodate}. Aujourd'hui, plus de 200 trésors et 230 000 monnaies sont publiés. 
Beaucoup de tables composent cette base, mais les principales sont les tables \enquote{Ensemble}, qui décrit les monnaies provenant de trésors, monnaies d'or isolées ou dépôts monétaires du territoire national, mais aussi les monnaies issues de fouilles, \enquote{monnaies} et \enquote{objets}. Ces tables piochent leurs informations dans des tables de listes fermées, ce qui évite les doublons. Dans la base, les ensembles sont notamment visibles sous la forme d'une frise chronologique. 

\subsection{Répertoire des mains musicales}

Le projet est mené par le département de la Musique de la BnF depuis 2020, et vise à faciliter l'identification des écritures musicales, utilisant des technologies de data mining et de reconnaissance de formes. La base de données, qui est le premier volet, s'organise autour de trois entités : les scripteurs, les manuscrits et les signes musicaux et compte actuellement 152 scripteurs, 359 manuscrits et 3454 signes musicaux. Le deuxième volet du projet consiste à utiliser l'intelligence artificielle pour développer un outil informatique capable de reconnaître les scripteurs de manière automatique\footcite{noauthor_repertoire_nodate}. 
 